\section{Introduction}
Lorem ipsum dolor sit amet, consectetur adipiscing elit. Sed do eiusmod tempor incididunt ut labore et dolore magna aliqua. Ut enim ad minim veniam, quis nostrud exercitation ullamco laboris nisi ut aliquip ex ea commodo consequat. Duis aute irure dolor in reprehenderit in voluptate velit esse cillum dolore eu fugiat nulla pariatur. Excepteur sint occaecat cupidatat non proident, sunt in culpa qui officia deserunt mollit anim id est laborum.

\subsection{State of the art}
Sed ut perspiciatis unde omnis iste natus error sit voluptatem accusantium doloremque laudantium, totam rem aperiam, eaque ipsa quae ab illo inventore veritatis et quasi architecto beatae vitae dicta sunt explicabo. Nemo enim ipsam voluptatem quia voluptas sit aspernatur aut odit aut fugit, sed quia consequuntur magni dolores eos qui ratione voluptatem sequi nesciunt. Neque porro quisquam est, qui dolorem ipsum quia dolor sit amet, consectetur, adipisci velit, sed quia non numquam eius mo....

\paragraph{Linkability}

The EUDI-ARF standard dictates that credentials \textit{issued} to a
holder, can be \textit{presented} (in the form of a verifiable
presentation) to a relying party in order to have one or more
attributes verified. Every verifiable presentation includes one or
more HMAC(s), formatted in SD-JWT: the HMACs are identical each time a
verifiable presentation is produced from a certain credential. This
makes possible for colluding relying parties, or to malicious actors,
to trace a holder's identity by collecting, exchanging and confronting
verifiable presentations (linkability). This threat appears to be well
mitigated by BBS+ through its Zero Knowledge Proof implementation.

\section{Overview}

\subsection{Feature Comparison}

\begin{table}[]
\centering
\renewcommand{\arraystretch}{2}
\begin{tabular}{ccccc}
\hline
       & \textbf{UP} & \textbf{UR} & \textbf{TR} & \textbf{URG} \\ \hline
SD-BLS & no          & yes         & yes         & yes \\ \hline
SD-JWT & no          & wip         & no          & no  \\ \hline
BBS+   & yes         & wip         & no          & yes \\ \hline
\end{tabular}
\vspace*{5mm}
\caption{feature comparison}
\label{tab:features}
\end{table}

We briefly round up on feature differences between the named selective
disclosure cryptographic schemes, as shown in table
\ref{tab:features}, mainly distinguishing between four fundamental
features:
\begin{itemize}
\item UP: Unlinkable Presentation
\item UR: Unlinkable Revocation
\item TR: Threshold Revocation
\item URG: Unregroupability
\end{itemize}

Where \emph{wip} is mentioned, it means work in progress on adoption of
bitstring status lists for unlinkable revocations.

The ``unregroupability'' feature refers to the fact that different
presentations of different claims cannot be linked to each other
(regrouped) as presented by the same holder, even when they are signed
by the same issuer.

\subsection{Key contributions}
The cryptographic scheme described in this paper, named
\textit{SD-BLS} for brevity, implements all the properties of the
SD-JWT scheme and proposes a novel cryptographic approach to similar
data structures. Furthermore SD-BLS proposes novel anonymous
cryptographic revocation flow for verifiable credentials, that aims to
solve governance issues posed by status and revocation lists.

\paragraph{Selective disclosure}
Similarly to the SD-JWT and mDOC formats, SD-BLS produces an array of
claims: the elements of the array are individually signed by the
issuer. In SD-BLS the signature(s) replace the HMAC and still enable
the holder to selectively disclose only certain signed credentials,
and produce a \textit{proof of possession} that minimizes private
information given to verifiers.

\paragraph{Anonymous cryptographic revocation}
SD-BLS proposes a novel approach to credential revocation: the data
published by the revocation issuer will produce cryptographic material
that contains no information about the credential holders. The
cryptographic revocation material allows anyone to verify if an SD-BLS
proof produced by a credential holder has been revoked. The
unlinkability, and thus anonymity of the cryptographic revocation,
allows the revocation issuer to share revocations in public and allows
anyone to verify if credentials have been revoked.

\paragraph{Multi-stakeholder governance of revocations}
SD-BLS allows to introduce a new trusted party to the issuance phase:
the revocation dealer. The dealer doesn't need to know the content of
the credential or the identity of an holder: its role is that of
producing the revocation signature and distributing its secret key to
a configurable range of revocation issuers. Later on a configurable
quorum of issuers may reconstruct the revocation secret key to revoke
a credential. This way the decision on a revocation is not delegated
to a single peer, but to a multi-stakeholder governance that protects
the holder from issuer corruption.

\paragraph{Fast revocation checks}
SD-BLS shows good benchmark results on revocation checks: the main
operation a verifier needs to do is an additional $\mathbb{G}_2$
multiplication. The cost of this operation grows linearly with the
number of published revocations.

\subsection{Applications}
In this section we present some applications and use cases for digital
identity and credentials, that could benefit from using the SD-BLS
scheme.

\paragraph{Digital identity}
The focus of the EUDI-ARF specifications is identity documents: it
defines mechanisms and data structures to issue a Personal
Identification (PID) as, for instance, with digital driving
licenses. Similarly, the US government is experimenting with W3C-VC
and mDOC for cross-states interoperable driving licenses. The SD-BLS
data format is similar to SD-JWT and mDOC, offering selective
disclosure and anonymous revocation.

\paragraph{Academic credentials}
Diploma and academic credentials are among the core offerings of EBSI
as well as a primary research target of the W3C VC working group.

\paragraph{KYC/AML}
We are unaware of standardization efforts for interoperable
credentials in the fields of "Know Your Customer" (KYC) and Anti
Money-Laundering (AML) certifications. We are aware of solution
providers experimenting with W3C-VC for AML applications and believe
SD-BLS can greatly improve the governance of credential revocation,
which is a critical component for this use case.

\paragraph{Generic \textit{light} credentials}
As the digital identity and verifiable credential technologies are
maturing, they are being considered for usage in less privacy
concerning applications, such subscriptions and membership and
fidelity cards.

\paragraph{Verifiable credentials on Blockchain}
SD-BLS can be used with blockchain-based smart-contracts to activate
certain functions:
\begin{itemize}
    \item A proof of possession can be published and be peer-verified
      on-chain in its private form without disclosing the credential
      contents.
    \item A smart-contract may verify if one or more holder's
      credentials match the requirement needed to process a
      transaction.
    \item Issuers can publish their cryptographic revocation lists on
      chain, allowing smart-contracts to verify the status of a
      credential.
    \item The issuer's public keys can also be published on-chain,
      although this does not represent a novelty.
\end{itemize}

\section{Implementation}

In this section we will provide a detailed description of the
algorithm we propose for selective disclosure and unlinkable
revocation using BLS signatures.

\subsection{Notations and assumptions}

We will adopt the following notations:
\begin{itemize}

\item $\mathbb{F}_p$ is the prime finite field with $p$ elements
  (i.e. of prime order $p$);

\item $E$ denotes the (additive) group of points of the curve
  BLS12-381 \cite{bls381-12} which can be described with the
  Weierstrass form $y^2=x^3 + 16$;

\item $E_T$ represents instead the group of points of the twisted
  curve of BLS12-381, with embedding degree $k=12$. The order of
  this group is the same of that of $E$;

\end{itemize}

We also require the notion of a cryptographic pairing. \cite{weilpair}

For the purpose of our protocol we will consider the \emph{Miller
pairing} $e: E_T \times E\to \mathbb{G}_T$, where $\mathbb{G}_T\subset
\mathbb{F}_{p^{12}}$ is the subgroup containing the $n$-th roots of
unity, and $n$ is the order of the groups $E$ and $E_T$.\\ For
completeness we also recall the main properties of the map:

\begin{itemize}

\item [i.] \emph{Bilinearity}, i.e. given $P_1,Q_1\in E_T$
  and $P_2,Q_2\in E$, we have
  \begin{align*}
    e(P_2,P_1+Q_1) = e(P_2,P_1)\cdot e(P_2,Q_1)   \\
    e(P_2+Q_2,P_1) = e(P_2,P_1)\cdot e(Q_2,P_1)
  \end{align*}

\item[ii.] \emph{Non-degeneracy}, meaning that for all
  $g_1\in E_T, g_2\in E$, $e(g_2,g_1)\ne
  1_{\mathbb{G}_T}$, the identity element of the group
  $\mathbb{G}_T$;

\item[iii.] \emph{ Efficiency}, so that the map $e$ is easy to
  compute;

\item[iv. ] $E_T \ne E$, and moreover, that
  there exist no efficient homomorphism between $E_T$ and
  $E$.

\end{itemize}


\section{Acknowledgements}
At vero eos et accusamus et iusto odio dignissimos ducimus qui blanditiis praesentium voluptatum deleniti atque corrupti quos dolores et quas molestias excepturi sint occaecati cupiditate non provident, similique sunt in culpa qui officia deserunt mollitia animi, id est laborum et dolorum fuga. Et harum quidem rerum facilis est et expedita distinctio. Nam libero tempore, cum soluta nobis est eligendi optio cumque nihil impedit quo minus id quod maxime placeat facere possimus, omnis voluptas assumenda est, omnis dolor repellendus.

\bibliographystyle{IEEEtran}
\bibliography{references}

\end{document}
